\chapter*{Statement on the use of AI tools}
\addcontentsline{toc}{chapter}{Statement on the use of AI tools}

A large-language-model tool (Claude, Anthropic; successive 2025--2026 model releases) was used
during this work, through an agentic command-line interface that can read and edit files, run
programs, and inspect their output. In the interest of transparency, this statement records what the tool was used for,
what it was not used for, and how its output was controlled.

\paragraph{What the tool was used for.}
\begin{itemize}
  \item \textbf{Implementation assistance.} Drafting and revising code under my direction,
  in particular on the reference side of the toolchain---the Dr.Jit array-programming idiom
  and the reference's actively developed research codebases, which span two
  repositories---and for repetitive engineering work (build plumbing, refactors).
  \item \textbf{Experiment orchestration.} Executing experiments of my design: writing
  batch scripts around the renderers, launching and monitoring long render campaigns,
  collecting outputs, and tabulating results.
  \item \textbf{Analysis and figures.} Producing plotting and statistics scripts
  (matplotlib/numpy) for the measurements reported here, under conventions I specified.
  \item \textbf{Verification tooling.} Implementing the deterministic probe scripts used in
  the validation and correction work (e.g.\ the segment-additivity and quadrature checks of
  \Cref{sec:results-firefly} and \Cref{app:corrections}), whose outputs are reproducible by
  running the scripts against the relevant code.
  \item \textbf{Prose.} Drafting and editing passes over text I outlined, and improving the
  readability of text I wrote.
\end{itemize}

\paragraph{What was not delegated.} The research direction; the renderer's architecture and
every design decision in it; the choice, design, and interpretation of every experiment; the
decisions about what to claim and how to qualify it; and the final wording. I have reviewed
the material in this document and take full responsibility for every claim in it.

\paragraph{How the output was controlled.} AI-produced code and analysis were treated as
untrusted input, subject to the same standard of evidence as any third-party code. The same
standard shapes \Cref{ch:validation}: the central claims of this thesis are supported by
checks that can be re-run ---
closed-form and energy-conservation tests with exactly known answers, deterministic
replays, cross-renderer agreement against an independent implementation, and archived render
artifacts from which every reported number reproduces via scripts kept alongside the
sources. Where AI assistance materially contributed to a specific technical finding, the
verification is deterministic and independent of the tool that helped find it. The clearest
instance is the set of reference corrections of \Cref{sec:results-firefly}: each of the
five findings is certified by a parameter-free probe that any reader can re-run.
